\chapter{Nature of the Temporal State}

The temporal architecture developed in this work is organized around an ordered relational support, a state carried on that support, directed event activity, transport between event surfaces, local bifurcation at admitted events, an internal elapsed measure, and retained memory lineage. These objects are typed separately and are joined by the dependency structure developed in the preceding and subsequent formal layers.

\section{Ordered relational support}

Let
\[
(S,\prec)
\]
be an ordered relational domain and
\[
\boxed{\Psi:S\to\mathcal H}
\]
the temporal-state map. The relation $\prec$ supplies precedence between admitted events, while $\Psi$ supplies the state assigned to the ordered support.

For an admitted transition $a\to b$, the directed kinetic representation is
\[
W_{a\to b}=M_{ab}e^{A_{ab}/2},
\qquad
W_{b\to a}=M_{ab}e^{-A_{ab}/2}.
\]
The associated positive activity and directed current are
\[
\boxed{\mathfrak a_{ab}=2M_{ab}\cosh(A_{ab}/2)>0,}
\]
\[
\boxed{\mathfrak j_{ab}=2M_{ab}\sinh(A_{ab}/2).}
\]
They separate elapsed pace from transition orientation. Their ratio reconstructs the directed affinity,
\[
\boxed{
A_{ab}=2\operatorname{artanh}\!\left(\frac{\mathfrak j_{ab}}{\mathfrak a_{ab}}\right).
}
\]

\section{NOW, bifurcation and transport}

The NOW layer localizes admitted event support. At an event $s_n$, the state update is represented by
\[
\boxed{
\Psi_T(s_n^+)=B_n\Psi_T(s_n^-).
}
\]
For the admitted reversible reference family,
\[
B(q)=e^{qG},
\]
with the unitary directional subclass
\[
B_\phi(\beta)=e^{-i\beta G}.
\]

Between admitted event surfaces, local transport is represented by
\[
\boxed{\Psi_{n+1}=U_n\Psi_n.}
\]
An ordered history therefore has the compositional form
\[
\boxed{
\Psi_N
=U_{N-1}B_{N-1}\cdots U_2B_2U_1B_1\Psi_0,
}
\]
with operator ordering inherited from $(S,\prec)$.

\section{Internal elapsed activity}

The positive temporal activity supplies the internal elapsed one-form
\[
\boxed{
d\tau_{\rm int}=\phi\,d\lambda,
\qquad
\phi=\frac{\mathfrak a}{\mathfrak a_\star}>0.
}
\]
Here $\lambda$ labels the ordered relational path and $d\tau_{\rm int}$ measures accumulated elapsed activity along it.

For a local relational subsystem and an admitted reference clock,
\[
d\tau_x=\phi_xd\lambda,
\qquad
d\tau_{\rm ref}=\phi_{\rm ref}d\lambda,
\]
the relational lapse is
\[
\boxed{
N_R(x)=\frac{d\tau_x}{d\tau_{\rm ref}}
=\frac{\phi_x}{\phi_{\rm ref}}>0.
}
\]
The ratio is invariant under the admitted increasing reparameterization of $\lambda$ and composes under a change of reference clock,
\[
N_{x|s}=N_{x|r}N_{r|s}.
\]
After explicit reference-clock calibration,
\[
dt=T_{\rm ref}d\tau_{\rm ref},
\]
the local calibrated elapsed interval becomes
\[
\boxed{d\hat\tau_x=N_Rdt.}
\]

On the admitted lapse/phase-rate bridge,
\[
\boxed{r_t=N_Rr_n^{(\tau)},}
\]
so the local phase rate and the relational elapsed-clock ratio share the same temporal normalization.

\section{Memory as temporal lineage}

The Memory layer retains the consequences of admitted events in ordered internal elapsed activity. ORCHORBITAL organizes the resulting history through the active-attractor sequence
\[
\boxed{\alpha=(a_1,\ldots,a_N)}
\]
and ordered signed winding
\[
\boxed{\mathcal W=(\Delta W_1,\ldots,\Delta W_N).}
\]
The residence-bound radial lineage adds
\[
\boxed{
\mathcal R=(\rho_1,\ldots,\rho_N),
\qquad
\rho_k=\|r_k-c_{a_k}\|>0,
}
\]
with every radial coordinate bound to its exact event-residence cell and committed post-segment state.

The retained temporal state therefore contains both the current state and an ordered lineage used by the declared reconstruction gates.

\section{Retrodiction and hidden history geometry}

Within one retained active-attractor stratum, the winding--radius carrier supplies the constructive lift
\[
(\alpha,\mathcal W,\rho_1,\ldots,\rho_{N-1},r_N)
\longmapsto
(r_1,\ldots,r_N),
\]
which composes with the exact position-lineage inverse
\[
(r_1,\ldots,r_N)
\longmapsto
(u_1,\ldots,u_N).
\]
Spatial Offset Divergence audits collision fibers of compressed observation maps by measuring the separation between the full position histories carried by observation-equivalent candidates.

This links the temporal-state description to a concrete reconstruction question: how much ordered lineage must be retained for an observation to select one history in the declared domain.

\section{Typed temporal architecture}

The resulting structure is summarized by
\[
\boxed{
\text{ORDER}
\to
\text{EVENT SUPPORT}
\to
\text{BIFURCATION}
\to
\text{TRANSPORT}
\to
\text{ELAPSED ACTIVITY}
\to
\text{MEMORY LINEAGE}
\to
\text{RETRODICTION}.
}
\]
The mathematical roles are correspondingly separated into precedence, state, activity, directed current, bifurcation, transport, elapsed one-form, relational lapse, retained lineage and inverse-history map. Physical clock calibration and downstream spacetime closure retain their declared evidence and dependency gates.
